\chapter{Limites de l'application actuelle}

L'analyse des pratiques manuelles actuelles et des tentatives d'automatisation locale met en évidence six limites majeures :

\begin{enumerate}
    \item \textbf{Absence de filtrage et de conversion automatisés :} Obligation d'effectuer des copier-coller manuels répétitifs et de traiter chaque feuille de calcul séparément.
    \item \textbf{Dégradation de la mise en page et manque de standardisation :} Débordements fréquents des tableaux hors des marges Word et hétérogénéité des styles d'affichage appliqués par les agents.
    \item \textbf{Dépendance forte aux postes clients :} Nécessité d'installer Microsoft Office et des scripts VBA locaux sur chaque ordinateur, nuisant à la portabilité multiplateforme.
    \item \textbf{Lourdeur opérationnelle et perte de productivité :} Mobilisation chronophage des agents (plusieurs heures pour un seul rapport) sur des tâches répétitives sans valeur ajoutée.
    \item \textbf{Risques majeurs de sécurité et fuites de données :} Recours non contrôlé à des services cloud publics gratuits exposant des documents confidentiels, et vulnérabilités associées aux macros VBA executables.
    \item \textbf{Déficit d'audit et de gouvernance des données :} Absence totale de registre centralisé, d'horodatage certifié et de vérification d'intégrité (empreinte SHA-256) des rapports produits.
\end{enumerate}

L'ensemble de ces limites justifie la mise en place immédiate d'une plateforme web centralisée, souveraine et sécurisée au sein de l'ANTIC.
